\chapter{Changes Across the Scientific Workflow}
\label{ch:scientificworkflow}

Language models can assist at several stages of research, but each use needs
different checks. Following the workflow from literature search to review
helps us see where they are useful and where judgement remains ours.

\section{Finding literature}

This stage has changed most. \textcite{zhang2025deepresearch} trace the path from
keyword search through RAG to \term{deep research}: systems that break a question
into parts, search, read what they find, and decide what to search next. For a
many-sided question --- thermal tolerance in Mediterranean seagrass, including
grey literature and non-English sources --- this beats manual search on both
coverage and time. In conservation, where systematic reviews are slow and
expensive, \textcite{bergertal2024evidence} conclude that judicious use of AI
could speed evidence synthesis up.

\begin{table}[htbp]
\centering
\caption{Literature sources, and what each is good for.}
\label{tab:literature}
\small
\begin{tabularx}{\textwidth}{lX}
\toprule
\textbf{Resource} & \textbf{Use} \\
\midrule
PubMed & Biomedical and molecular; excellent structured metadata \\
OpenAlex & Open bibliographic graph with a free API; good for building your own pipeline \\
Google Scholar & Broadest coverage including grey literature; no usable API \\
Biodiversity Heritage Library & Historical taxonomic literature; irreplaceable \\
Biodiversity PMC & Biodiversity-specific full text \parencite{pasche2026biodivpmc} \\
Web of Science, Scopus & Curated, subscription; citation analysis \\
Zenodo / BioSysLit & Preprints, datasets, community collections \\
\bottomrule
\end{tabularx}
\end{table}

Resolve every reference. About 147,000 invented citations appeared in 2025
alone \parencite{zhao2026citations}, each with real authors and a plausible
title. A short script catches them:

\begin{lstlisting}[style=shell,caption={Resolve every DOI before it enters a manuscript}]
grep -oP 'doi\s*=\s*\{\K[^}]+' references.bib | while read -r doi; do
    code=$(curl -s -o /dev/null -w "%{http_code}" "https://doi.org/$doi")
    [ "$code" = "200" ] || [ "$code" = "302" ] || echo "UNRESOLVED: $doi"
done
\end{lstlisting}

Resolving is not enough on its own: the returned title must match, and entries
without a DOI need checking by hand. This book's bibliography was checked that
way.

\begin{keypoint}
Never cite a paper you have not opened. An AI-assisted search finds candidates;
it does not replace reading them.
\end{keypoint}

\section{Reading and analysis}

With the PDF in front of it, a model can explain a method, translate, or answer
questions about the paper; asked from memory, it invents details. What it cannot
do is judge whether the paper is any good --- whether the sample is adequate, the
controls appropriate, the conclusion supported. A systematic review of 63
studies found that models agreed with human screeners about nine times in ten,
but their median accuracy in assessing risk of bias was 0.62
\parencite{laignelot2026systematicreviews}. A confident summary of a bad
paper reads exactly like one of a good paper. For analysis, use models to write
code, not to produce results (\Cref{ch:agentic,ch:reproducibility}).

\section{Writing}

These tools help with structure, first drafts and proofreading, and the authors
remain responsible for the content \parencite{salvagno2023scientificwriting,
granjeiro2025futurewriting}. The deeper issue is voice.
\textcite{alhosni2025preservingvoice} reviews evidence that generative tools make
academic prose more uniform, blurring what the author wrote.
\textcite{singh2025authorialvoice} describe a tool built to develop an author's
own voice rather than replace it. This distinction matters in practice.
When a tool helps us develop our own draft, writing remains part of the
thinking process. Delegating the draft can remove that opportunity to learn.

Journal policy has converged \parencite{naturemethods2026responsibly}.
AI cannot be an author, because authorship requires accountability. Authors
must disclose its use specifically and remain responsible for the content.
Uploading a manuscript under review to a hosted service breaches confidentiality. That
last point is often missed. A corpus-level analysis estimated that 6.5 to 16.9
per cent of the text in peer reviews for four AI conferences in 2023 and 2024
was substantially modified by a model \parencite{liang2024peerreview}. Claims
based on detectors applied to single reviews need more caution, because the
detectors also flag reviews that a model has only polished
\parencite{saha2026peerreview}. If you want help with a review, use a local model or
none --- one of the clearest cases for local-first in this book.

\section{Review, assessment and detection}

Automated checks --- resolving references, checking reported statistics,
spotting duplicated figures --- are legitimate and need no generative model.
Automated review \emph{judgement} is not: a model produces text shaped like a
review without the judgement. Assessments such as IPBES require every statement
to trace to a source and carry a calibrated confidence
\parencite{krug2026ipbes}; generative summaries break the first and invent the
second. Use models there to find literature, extract data into a consistent
schema, translate and flag inconsistencies --- not to write the assessment. And
do not use AI-text detectors to judge people (\Cref{ch:powers}): make
disclosure easy and normal instead.
